%% LyX 2.4.4 created this file.  For more info, see https://www.lyx.org/.
%% Do not edit unless you really know what you are doing.
\documentclass[english]{article}
\usepackage[T1]{fontenc}
\usepackage[latin9]{inputenc}
\usepackage{units}
\usepackage{enumitem}
\usepackage{amsmath}
\usepackage{amsthm}
\usepackage{geometry}
\geometry{verbose,tmargin=1in,bmargin=1in,lmargin=1in,rmargin=1in}

\makeatletter
%%%%%%%%%%%%%%%%%%%%%%%%%%%%%% Textclass specific LaTeX commands.
\numberwithin{equation}{section}
\numberwithin{figure}{section}
\newlength{\lyxlabelwidth}      % auxiliary length 

%%%%%%%%%%%%%%%%%%%%%%%%%%%%%% User specified LaTeX commands.
\usepackage{fancyhdr}
\usepackage{lastpage}
\pagestyle{fancy}
\usepackage{enumitem}
\usepackage{ifthen}
\everymath=\expandafter{\the\everymath\displaystyle}
%\DeclareMathSizes{10}{10}{10}{10}
\usepackage{multicol}

\usepackage{tikz}
\usepackage{tikz-3dplot}
\usetikzlibrary{calc,arrows}

\usetikzlibrary{patterns}
\usetikzlibrary{plotmarks}
\usepackage{pgfplots}
\pgfplotsset{compat=newest}

\usepackage{calc}
\usepackage[nomessages]{fp}

\usepackage{datetime}
% Added by lyx2lyx
\setlength{\parskip}{\smallskipamount}
\setlength{\parindent}{0pt}

\makeatother

\usepackage{babel}
\begin{document}
\renewcommand{\author}{Dr.\,James Ashe}

\newcommand{\classNum}{336}

\newcommand{\classSec}{A}

\newcommand{\examType}{HW 5.2}

\renewcommand{\date}{\formatdate{6}{11}{2013}}

%%First page's header%%%%%%%%%%%%%%%%%%%%%%
\fancypagestyle{firststyle} %%header settings for first pagr
{
	\fancyhead[L]{MTH \classNum{}(\classSec{}) \author{}\\
	\textbf{\examType{}} ~\date{}}
	\fancyhead[C]{}
	\fancyhead[R]{\textbf{Solutions}}
	\setlength{\headsep}{14pt}
}
%%%%%%%%%%%%%%%%%%%%%%%%%%%%%%%%%%%%%%%%%%

%%Next page's header%%%%%%%%%%%%%%%%%%%%%%%
\setlist{nolistsep}
\lhead{\classNum{}(\classSec{})} 
\chead{\textbf{Solutions}} 
\cfoot{\thepage{}~of~\pageref{LastPage}} 
\setlength{\headsep}{5pt} 
%%%%%%%%%%%%%%%%%%%%%%%%%%%%%%%%%%%%%%%%%%%

%%Various settings; see the preamble for other settings%%%%%
%\setenumerate[0]{leftmargin=12pt,itemindent=0pt,labelwidth=0pt}
\setlist{topsep=.5pt}
\renewcommand{\labelenumi}{\textbf{\arabic{enumi}.}}
\renewcommand{\labelenumii}{\textbf{\alph{enumii}.}}
\thispagestyle{firststyle}
%%%%%%%%%%%%%%%%%%%%%%%%%%%%%%%%%%%%%%%%%%

\newcommand{\xmax}{0}
\newcommand{\xmin}{-\xmax}
\newcommand{\xstep}{0}
\newcommand{\ymax}{0}
\newcommand{\ymin}{-\ymax}
\newcommand{\ystep}{0} 
\newcommand{\dspace}{\vspace{1cm}}
\newcommand{\xa}{0}
\newcommand{\xb}{0}
\newcommand{\ya}{1}
\newcommand{\yb}{1}
\begin{enumerate}[resume]
\item \vspace{0cm}Note: these are the same matrices from 5.1 exercise 1.

\begin{enumerate}[resume]
\item Using the first row ($i=1$). So the signs are $+,-,+$\\
$\det\begin{bmatrix}\begin{array}{rrr}
-1 & 3 & 5\\
6 & 4 & 2\\
-2 & 5 & 1
\end{array}\end{bmatrix}=(-1)\det\begin{bmatrix}\begin{array}{rr}
4 & 2\\
5 & 1
\end{array}\end{bmatrix}-(3)\det\begin{bmatrix}\begin{array}{rr}
6 & 2\\
-2 & 1
\end{array}\end{bmatrix}+5\det\begin{bmatrix}\begin{array}{rr}
6 & 4\\
-2 & 5
\end{array}\end{bmatrix}=(-1)(4-10)-3(6+4)+5(30+8)=166$
\item Above and in many examples we used the first row, but you can use
any row (or column) you like. Since there are a row of zeros on the
fourth row $(i=4)$, we'll us that row. Giving the signs $-,+,-,+$\\
\begin{eqnarray*}
\det\begin{bmatrix}\begin{array}{rrrr}
1 & -1 & 0 & 1\\
0 & 2 & 1 & 1\\
2 & -2 & 2 & 3\\
0 & 0 & 6 & 2
\end{array}\end{bmatrix} & = & -(0)\det\begin{bmatrix}\begin{array}{rrr}
-1 & 0 & 1\\
2 & 1 & 1\\
-2 & 2 & 3
\end{array}\end{bmatrix}+(0)\det\begin{bmatrix}\begin{array}{rrr}
1 & 0 & 1\\
0 & 1 & 1\\
2 & 2 & 3
\end{array}\end{bmatrix}-(6)\det\begin{bmatrix}\begin{array}{rrr}
1 & -1 & 1\\
0 & 2 & 1\\
2 & -2 & 3
\end{array}\end{bmatrix}+(2)\det\begin{bmatrix}\begin{array}{rrr}
1 & -1 & 0\\
0 & 2 & 1\\
2 & -2 & 2
\end{array}\end{bmatrix}\\
 & = & 0+0-6(2)+2(4)\\
 & = & -4
\end{eqnarray*}
\end{enumerate}
\item \vspace{0cm}

\begin{enumerate}
\item $\det A=-2\det\begin{bmatrix}\begin{array}{rr}
2 & 1\\
4 & 2
\end{array}\end{bmatrix}+3\det\begin{bmatrix}\begin{array}{rr}
1 & 1\\
1 & 2
\end{array}\end{bmatrix}-0=3$ and $\det B_{2}=\det\begin{bmatrix}\begin{array}{rrr}
1 & 1 & 1\\
2 & 2 & 0\\
1 & -1 & 2
\end{array}\end{bmatrix}=-2\det\begin{bmatrix}\begin{array}{rr}
1 & 1\\
-1 & 2
\end{array}\end{bmatrix}+2\det\begin{bmatrix}\begin{array}{rr}
1 & 1\\
1 & 2
\end{array}\end{bmatrix}-0=-4$. So by Cramer's rule, $x_{2}=\frac{\det B_{2}}{\det A}=-\frac{4}{3}$.
\item By proposition 2.4, $A^{-1}=\frac{1}{\det A}C^{\mbox{T}}$, where
$C^{\mbox{T}}$ is the \emph{transpose }of the $C=\left[C_{ij}\right]$,
the matrix composed of $A$'s cofactors. The cofactor matrix is, $C=\begin{bmatrix}\begin{array}{rrr}
6 & -4 & 5\\
0 & 1 & -2\\
-3 & 2 & -1
\end{array}\end{bmatrix}.$ The transpose is $C^{\mbox{T}}=\begin{bmatrix}\begin{array}{rrr}
6 & 0 & -3\\
-4 & 1 & 2\\
5 & -2 & -1
\end{array}\end{bmatrix}$. So then $A^{-1}=\frac{1}{3}\begin{bmatrix}\begin{array}{rrr}
6 & 0 & -3\\
-4 & 1 & 2\\
5 & -2 & -1
\end{array}\end{bmatrix}=\begin{bmatrix}\begin{array}{rrr}
2 & 0 & -1\\
-\nicefrac{4}{3} & \nicefrac{1}{3} & \nicefrac{2}{3}\\
\nicefrac{5}{3} & -\nicefrac{2}{3} & -\nicefrac{1}{3}
\end{array}\end{bmatrix}$
\end{enumerate}
\item [\textbf{5.}]\vspace{0cm}

\begin{enumerate}
\item $\det\left(\begin{bmatrix}\begin{array}{rr}
1 & 5\\
2 & 4
\end{array}\end{bmatrix}-\begin{bmatrix}\begin{array}{rr}
t & 0\\
0 & t
\end{array}\end{bmatrix}\right)=\det\begin{bmatrix}\begin{array}{rr}
1-t & 5\\
2 & 4-t
\end{array}\end{bmatrix}=(1-t)(4-t)-(5)(2)=t^{2}-5t-6$
\item [\textbf{(e.)}]$\det\begin{bmatrix}\begin{array}{rr}
1-t & 1\\
-1 & 3-t
\end{array}\end{bmatrix}=(1-t)(3-t)-(-1)(1)=t^{2}-4t+4$
\item [\textbf{(f.)}]
\begin{eqnarray*}
\det\begin{bmatrix}\begin{array}{rrr}
-1-t & 1 & 2\\
1 & 2-t & 1\\
2 & 1 & -1-t
\end{array}\end{bmatrix} & = & (-1-t)\det\begin{bmatrix}\begin{array}{rr}
2-t & 1\\
1 & -1-t
\end{array}\end{bmatrix}-(1)\det\begin{bmatrix}\begin{array}{rr}
1 & 1\\
2 & -1-t
\end{array}\end{bmatrix}+(2)\det\begin{bmatrix}\begin{array}{rr}
1 & 2-t\\
2 & 1
\end{array}\end{bmatrix}\\
 & = & (-1-t)(t^{2}-t-3)-(-t-3)+2(2t-3)\\
 & = & -t^{3}+9t
\end{eqnarray*}
\end{enumerate}
\item [\textbf{7.}]\vspace{0cm}

\begin{enumerate}
\item By assumption, $\det A=\pm1$; so the cofactors of $A$ are $C_{ij}=(-1)^{i+j}\det(M_{ij})$,
which are all integers by exercise $6$ (proved in class). Thus $C=\left[C_{ij}\right]$
has all integer entries which implies that $C^{\mbox{T}}$ has all
integer entries. $A$ is nonsingular (since $\det A\neq0$), so by
proposition 2.4 $A^{-1}=\frac{1}{\det A}C^{\mbox{T}}=\pm C^{\mbox{T}}$.
Thus $A^{-1}$ also has all integer entries. 
\end{enumerate}
\end{enumerate}

\end{document}
